% =====================================================
% Plantilla adaptada de FloripaSat-2 / GOLDS-UFSC §4.5 + Apéndice A
% Para usar en: 00_documento_principal/chapters/02_sistema.tex (§2.2.4)
% =====================================================

\subsection{Cálculo de Enlace (Link Budget)}

El presupuesto de enlace se calcula para los enlaces de subida (uplink) y bajada
(downlink) en banda UHF/VHF, considerando el peor caso orbital: satélite en el horizonte
($\theta = 0^\circ$ de elevación). Los pasos detallados están en el Anexo de cálculo de enlace.

\begin{table}[h]
\centering
\caption{Resultados del link budget de INTISAT.}
\label{tab:link_budget}
\renewcommand{\arraystretch}{1.15}
\begin{tabular}{lrrl}
\toprule
\textbf{Parámetro}                       & \textbf{Uplink} & \textbf{Downlink} & \textbf{Unidad} \\
\midrule
Frecuencia                                & TBD   & TBD   & MHz \\
Potencia de transmisión                   & TBD   & TBD   & dBm \\
Ganancia de antena TX                     & TBD   & TBD   & dBi \\
Pérdida por cable + conectores            & TBD   & TBD   & dB \\
EIRP                                      & TBD   & TBD   & dBm \\
Pérdida en espacio libre (FSPL)           & TBD   & TBD   & dB \\
Pérdida por polarización                  & TBD   & TBD   & dB \\
Pérdida atmosférica + lluvia              & TBD   & TBD   & dB \\
Ganancia de antena RX                     & TBD   & TBD   & dBi \\
Temperatura de ruido del sistema          & TBD   & TBD   & K \\
G/T                                       & TBD   & TBD   & dB/K \\
\textbf{C/N$_0$ recibido}                 & TBD   & TBD   & dB·Hz \\
Bit rate                                  & TBD   & TBD   & bps \\
E$_b$/N$_0$ requerido (BER $10^{-5}$)     & TBD   & TBD   & dB \\
E$_b$/N$_0$ disponible                    & TBD   & TBD   & dB \\
\textbf{Margen}                           & \textbf{TBD} & \textbf{TBD} & \textbf{dB} \\
\bottomrule
\end{tabular}
\end{table}

\paragraph{Criterio de aceptación.}
El margen debe ser $\ge 3$~dB para considerar el enlace operativo. FloripaSat-2 logró
márgenes de $\ge 9{,}57$~dB en downlink y $\ge 17{,}81$~dB en uplink (con FEC).

\paragraph{Análisis de pasadas sobre estación INRAS.}
Considerando órbita TLE-objetivo y estación INRAS (PUCP, Lima), se espera:
\begin{itemize}
    \item Periodo de contacto máximo: TBD~s
    \item Periodo de contacto promedio: TBD~s
    \item Periodo de contacto total (60 días): TBD~s
    \item Transferencia teórica de datos: TBD~bits/día
\end{itemize}
